In this work we proposed a graph-guided prompting workflow intended to improve the reliability, traceability, and citation quality of long-form, section-level scientific writing.
The approach models a draft document as a directed document graph whose nodes correspond to sections and whose edges encode hierarchical and evidentiary dependencies.
Section-level generators are required to emit strict, machine-parseable JSON outputs that include structured text, explicit citation slots, and provenance metadata.
Core engineering components include a dependency-respecting section scheduler, per-node evidence retrieval with prompt injection, constrained generation that enforces citation and provenance fields for leaf sections, and an assembly stage that materializes a LaTeX document, refs.bib, and an audit trail linking claims to sources.
These elements are designed to bias generation toward conservative claims and to produce outputs amenable to automated validation and human review.
We emphasize three practical advantages: schema-constrained outputs enable automated validation and surfacing of missing or malformed citation metadata; the document-graph clarifies local context and generation order, improving targeted retrieval; and assembly-stage provenance yields machine-actionable artifacts for downstream verification.
We also document key limitations and caveats.
The pipeline inherits retrieval limitations: when relevant sources are absent or ranking is poor, constrained prompting cannot manufacture grounded citations and factual coverage will suffer \cite{web_deng_2025_the_next_phase_of_scientific_fact_checking_advan}.
The present evaluation demonstrates feasibility rather than domain-general guarantees, and larger cross-domain studies are needed to quantify robustness and failure modes.
We did not perform exhaustive online related-work retrieval during preparation, so claims about absolute novelty are intentionally withheld.
Future work includes scaling and diversifying retrieval sources and evaluating how recall and ranking affect citation coverage \cite{web_deng_2025_the_next_phase_of_scientific_fact_checking_advan}, integrating automated fact-checking and contradiction detection into assembly and validation, and broadening empirical evaluation with standardized benchmarks.
We also recommend improving human--AI interfaces to let authors correct retrieval results or alter the section graph, and pursuing tighter integration with verification tooling and reference managers so the audit trail becomes actionable in editorial workflows.